\chapter{Tendencias web actuales}
\label{cap:tendencias-web}

% TODO: integrar el aporte validado de Fajardo Montes Michael Xavier
% (investigacion sobre Jamstack, PWA e inteligencia artificial generativa
% aplicada al desarrollo web, con reflexion critica).
%
% Fuente prevista (no integrada todavia en este capitulo):
%   docs/u4/investigacion/TENDENCIAS-WEB.md
%
% Al integrar, trasladar aqui unicamente las fuentes que ese documento ya
% cite y que sean verificables (sin inventar autores, anios ni URLs), y
% anadirlas a referencias.bib de este informe.

\pendienteIntegracion{Investigaci\'on sobre tendencias web (Fajardo Montes Michael Xavier): Jamstack, Progressive Web Apps (PWA) e inteligencia artificial generativa aplicada al desarrollo web, con reflexi\'on cr\'itica sobre su aplicabilidad a BIOPET.}

\section{Alcance previsto}
\label{sec:tendencias-alcance}

Este cap\'itulo integrar\'a, cuando est\'e disponible, el an\'alisis de tres
tendencias del desarrollo web moderno:

\begin{itemize}[nosep]
  \item \textbf{Jamstack:} arquitectura basada en contenido pre-renderizado
    servido desde CDN, con l\'ogica din\'amica delegada a APIs.
  \item \textbf{Progressive Web Apps (PWA):} aplicaciones web instalables,
    con soporte offline mediante \emph{service workers}.
  \item \textbf{Inteligencia artificial generativa} aplicada al desarrollo
    web: usos actuales, riesgos y l\'imites documentados.
\end{itemize}

junto con una reflexi\'on cr\'itica sobre la aplicabilidad concreta de cada
tendencia al caso real de BIOPET (por ejemplo, si la naturaleza
transaccional y multi-rol del sistema favorece o no la adopci\'on de
Jamstack, o si el frontend Angular ya existente facilita convertir la
aplicaci\'on en una PWA).

No se redacta contenido definitivo en esta versi\'on del informe porque
depende de la investigaci\'on independiente asignada a Fajardo Montes
Michael Xavier, todav\'ia no incorporada a este documento.
